\subsection{Vechiul Testament}
% Adaugă aici 5-6 versete relevante din VT



\begin{displayquote}
    „A răspuns Samuel: «Au doară arderile de tot şi jertfele sunt tot aşa de plăcute Domnului, ca şi ascultarea glasului Domnului? Ascultarea este mai bună decât jertfa şi supunerea mai bună decât grăsimea berbecilor.»”
    (\textit{1 Regi} 15, 22)
\end{displayquote}

\begin{displayquote}
    „Şi se vor binecuvânta prin neamul tău toate popoarele pământului, pentru că ai ascultat glasul Meu.”
    (\textit{Facerea} 22, 18)
\end{displayquote}

\begin{displayquote}
    „Puţini din voi vor rămâne, deşi veţi fi fost ca stelele cerului, pentru că n-aţi ascultat glasul Domnului Dumnezeului vostru.”
    (\textit{Deuteronomul} 28, 62)
\end{displayquote}

\begin{displayquote}
    „Ascultă fiică şi vezi şi pleacă urechea ta şi uită poporul tău şi casa părintelui tău.”
    (\textit{Psalmul} 44, 12)
\end{displayquote}

\begin{displayquote}
    „Calea celui nebun este dreaptă în ochii lui, iar cel înţelept ascultă de sfat.”
    (\textit{Pildele lui Solomon} 12, 15)
\end{displayquote}

\begin{displayquote}
    „Domnul Dumnezeu Mi-a dat Mie limbă de ucenic, ca să ştiu să grăiesc celor deznădăjduiţi. În fiecare dimineaţă El deşteaptă, trezeşte urechea Mea, ca să ascult ca un ucenic.”
    (\textit{Isaia} 50, 4)
\end{displayquote}

\subsection{Noul Testament}
% Adaugă aici 1-2 versete relevante din NT (maxim 2 conform cerinței)

\begin{displayquote}
    „Şi noi ştim că Dumnezeu nu-i ascultă pe păcătoşi; dar de este cineva cinstitor de Dumnezeu şi face voia Lui, pe acesta îl ascultă.”
    (\textit{Ioan} 9, 31)
\end{displayquote}